\chapter{CONCLUSION}

\section{Summary of Findings}
This thesis developed and validated a multi-stage computer vision framework to mitigate pediatric undercounting in crowded clinics. The pipeline combines human allometry, margin-expanded crop classification, and temporal filtering to count pediatric patients accurately under steep CCTV camera geometries.

The main findings of this work are:
\begin{enumerate}[(i)]
    \item \textbf{Separable Convolution Geometry (Theoretical Evaluation):} Mathematical analysis of discrete image geometry showed that splitting the 2D Sobel gradient operator into rank-1 outer product vectors ($\mathbf{s}_y \otimes \mathbf{d}_x^T$) reduces operations from 9 per pixel to 6 per pixel. While not deployed in the final tracking pipeline, this $33.3\%$ arithmetic decrease provides a foundation for future low-power edge filters.
    \item \textbf{Distillation Data Bottleneck:} Tests showed that poor training data quality may have contributed to distillation performance drops. The lack of steep camera angle data and limited examples of dense crowding restricted the student models ($F_1 = 0.736$, $\text{MAPE} = 26.2\%$) in hospital scenes. These limitations prompted the shift to a multi-stage architecture.
    \item \textbf{Image Slicing and Dynamic Range:} Slicing Aided Hyper Inference (SAHI) with $640 \times 640$ tiles ($20\%$ overlap) prevented feature loss for distant children occupying fewer than $32 \times 32$ pixels. CIELAB CLAHE increased local contrast by $+34.2\%$ without altering chromaticity.
    \item \textbf{Preliminary Allometric Separation:} The skeletal cephalocaudal ratio $R_{\text{ceph}} = L_{\text{torso}}/L_{\text{leg}}$ produced a distinct biological gap ($\Delta R_{\min} = 0.204$). We observed this preliminary separation within our small evaluated sample ($n=6$). Pediatric values ($0.941\text{--}0.982$) remained clearly above adult values ($0.621\text{--}0.737$). Setting threshold $\tau = 0.80$ separated demographic classes effectively.
    \item \textbf{Tracking and Counting Reliability:} Combining a base spatial detector with a fine-tuned demographic classifier and weighted temporal smoothing suppressed single-frame false detections. Across the evaluated 700-frame CCTV benchmark corpus (45 ground-truth occupants), the multi-stage pipeline matched demographic counts precisely.
    \item \textbf{CPU Vectorization:} Hardware benchmarks showed that native FP32 ONNX execution with AVX2 SIMD instructions achieved $6.8\text{--}9.1$ FPS across CCTV resolutions ($110.2\text{--}146.6$ ms latency). This ran $3.98\times$ faster in pipeline latency than emulated FP16 ($438.3$ ms, $2.3$ FPS). Standard hospital CPUs can run the pipeline without discrete GPUs.
\end{enumerate}

\section{Conclusions}
The evidence supports pediatric detection feasibility in non-clinical stress videos. Hospital videos primarily validated the camera environment on adults, meaning clinical pediatric validation remains necessary. Vision models require mathematical and biological constraints derived from image formation and human body growth.

The framework treats images as quantized matrices and applies Complete IoU regression. It crops candidate targets for classification, anchors demographic identities, and stabilizes tracks over time. This approach provides an objective, privacy-preserving occupancy counting tool. When integrated into clinical workflows, the system has the potential to provide objective telemetry to support triage operations and LHIMS-linked planning, pending larger-scale clinical pediatric validation.

\section{Recommendations for Clinical Deployment}
To deploy this system in Ghanaian hospitals, we recommend four practical steps:
\begin{enumerate}[(i)]
    \item \textbf{Hardware Specifications:} Hospitals should deploy the model on standard edge computers with x86\_64 processors supporting AVX2 SIMD extensions. Running native FP32 ONNX models avoids the overhead of FP16 software emulation and delivers reliable 7--8 FPS edge tracking throughput at low cost.
    \item \textbf{Camera Mounting and Angles:} CCTV cameras in waiting areas should be mounted at ceiling heights of $2.8\text{--}3.5$ meters with a downward angle of $25^{\circ}\text{--}35^{\circ}$. This angle captures the full torso and leg segments ($L_{\text{torso}}$ and $L_{\text{leg}}$) while reducing perspective distortion.
    \item \textbf{Patient Privacy Protection:} IT staff must preserve the zero-storage privacy design. Frames must be processed in RAM and discarded immediately. The system should send only anonymized numerical data (timestamps, track IDs, coordinates, and counts) to LHIMS.
    \item \textbf{Integration with Triage Dashboards:} Census outputs should feed directly into emergency nursing screens. The system should alert staff when pediatric queue numbers exceed thresholds under the South African Triage Scale (SATS).
\end{enumerate}

\section{Limitations and Future Research}

\subsection{Study Limitations}
This study identifies four primary limitations:
\begin{enumerate}[(i)]
    \item \textbf{Dataset Scale and Sequence Correlation:} We collected and annotated real-world clinic training images and partitioned them by unique video sequence. However, expanding to multi-hospital datasets with sequence-independent splitting across diverse clinical demographics remains necessary to prevent visual overfitting to local clothing styles.
    \item \textbf{Sample Size of Allometric Observations:} The skeletal cephalocaudal ratio was evaluated on an empirical sample ($n=6$ subjects). Due to strict ethical limitations on measuring and recording children in clinical triage, we validated the biological gap on a small, fully consented stress-test cohort before attempting wider clinical deployment. While findings match developmental anthropology, measuring larger cohorts across varied postures remains future work.
    \item \textbf{Curated Evaluation Corpus Scope:} The $100.0\%$ headcount accuracy reflects consolidated tracklets on a curated 700-frame benchmark (45 occupants). Raw single-frame detection achieves $67.1\%$ mAP@50. Unconstrained hospital queues with extreme crowding will experience tracking fragmentations.
    \item \textbf{Monocular Depth and Tilt Ambiguity:} Single 2D CCTV cameras cannot measure 3D spatial depth directly. Non-vertical body tilt or bending introduces angular foreshortening that degrades 2D ratio estimates.
\end{enumerate}

\subsection{Future Research Directions}
We identify three directions for future research:
\begin{enumerate}[(i)]
    \item \textbf{Infrared Night Monitoring:} Under zero-lux night lighting, CCTV cameras switch to infrared illumination. This removes color information ($a^*, b^*$). Future work should examine multi-spectral and thermal camera integration.
    \item \textbf{Campus-Scale Multi-Camera Tracking:} The system implements a lightweight cross-camera tracking plugin using a shared gallery of upper and lower body color histograms. While this operates across adjacent rooms, it depends on consistent clothing appearance and camera color calibration. In large hospitals with varied lighting between halls, waiting corridors, and consultation rooms, appearance matching can degrade. Future work should develop spatio-temporal camera transition models and graph neural networks without biometric facial logging.
    \item \textbf{3D Depth Sensing:} Adding lightweight stereo or Time-of-Flight (ToF) depth sensors would provide 3D spatial points. This would help address vertical physical obstructions in dense queues directly through 3D geometry.
\end{enumerate}
